\documentclass[12pt,a4paper]{article}

\usepackage[utf8]{inputenc}
\usepackage[T1]{fontenc}
\usepackage[spanish]{babel}
\usepackage{amsmath,amssymb}
\usepackage{geometry}
\usepackage{booktabs}
\usepackage{hyperref}
\usepackage{listings}
\usepackage{xcolor}
\usepackage{parskip}
\usepackage{fancyhdr}
\usepackage{array}
\usepackage{mathtools}

\geometry{margin=2.5cm, headheight=15pt}

\definecolor{codebg}{RGB}{245,245,245}
\definecolor{keywordcolor}{RGB}{0,100,180}
\definecolor{commentcolor}{RGB}{80,110,80}

\lstset{
  language=Python,
  backgroundcolor=\color{codebg},
  frame=single, rulecolor=\color{gray!40},
  basicstyle=\ttfamily\small,
  keywordstyle=\color{keywordcolor}\bfseries,
  commentstyle=\color{commentcolor}\itshape,
  numbers=left, numberstyle=\tiny\color{gray}, numbersep=8pt,
  breaklines=true, showstringspaces=false, tabsize=4,
}

\pagestyle{fancy}
\fancyhf{}
\rhead{Modelado y Simulaci\'on --- 22/04/2026}
\lhead{Punto 1 --- Docente: Omar C\'aceres}
\cfoot{\thepage}

\title{\textbf{Examen --- Modelado y Simulaci\'on}\\[0.3em]
       {\large Punto 1: Ra\'ices de $f(x)=2x\cos x-(x-2)^2$}}
\author{Juan Bautista Espino}
\date{22 de abril de 2026}

\begin{document}
\maketitle\thispagestyle{fancy}

% ══════════════════════════════════════════════════════════════════════════════
\section*{Enunciado}
% ══════════════════════════════════════════════════════════════════════════════

Dada la funci\'on
\[
  f(x) = 2x\cos(x) - (x-2)^2
\]
demostrar mediante el Teorema de Bolzano que existe al menos una ra\'iz en
$[0.8,\,1.4]$ y luego aproximarla con:
\begin{itemize}
  \item[(a)] M\'etodo de Punto Fijo acelerado con Steffensen-Aitken, $x_0=1$, 6 cifras de precisi\'on.
  \item[(b)] M\'etodo de Newton-Raphson, $x_0=1$, tolerancia $<10^{-8}$. Comparar ambos m\'etodos.
\end{itemize}

% ══════════════════════════════════════════════════════════════════════════════
\section*{Paso previo: Teorema de Bolzano}
% ══════════════════════════════════════════════════════════════════════════════

$f$ es continua en $\mathbb{R}$ (composici\'on de funciones continuas). Evaluamos en
los extremos del intervalo:

\[
  f(0.8) = 2(0.8)\cos(0.8) - (0.8-2)^2
         = 1.6 \times 0.696707 - 1.44
         = 1.114731 - 1.44
         \approx -0.325\,269
\]
\[
  f(1.4) = 2(1.4)\cos(1.4) - (1.4-2)^2
         = 2.8 \times 0.169967 - 0.36
         = 0.475907 - 0.36
         \approx +0.115\,908
\]

Como $f(0.8) < 0$ y $f(1.4) > 0$, se cumple:
\[
  f(0.8)\cdot f(1.4) \approx -0.037\,701 < 0
\]

Por el \textbf{Teorema de Bolzano} existe al menos un $c\in(0.8,\,1.4)$ tal que
$f(c)=0$. \hfill$\blacksquare$

% ══════════════════════════════════════════════════════════════════════════════
\section*{Parte (a) --- Steffensen-Aitken}
% ══════════════════════════════════════════════════════════════════════════════

\subsection*{Construcci\'on de $g(x)$}

Partimos de $f(x)=0$:
\[
  2x\cos(x) = (x-2)^2 \implies x-2 = -\sqrt{2x\cos x}
  \quad\text{(negativo pues $x\approx 1 < 2$)}
\]
\[
  \boxed{g(x) = 2 - \sqrt{2x\cos x}}
\]

\subsection*{Condici\'on de Lipschitz}

Derivando:
\[
  g'(x) = -\frac{\cos x - x\sin x}{\sqrt{2x\cos x}}
\]

En $x=1$:
\[
  g'(1) = -\frac{\cos 1 - \sin 1}{\sqrt{2\cos 1}}
        = -\frac{0.5403 - 0.8415}{\sqrt{1.0806}}
        = \frac{0.3012}{1.0395}
        \approx 0.2898
\]

Si bien $|g'(1)|=0.29 < 1$ en la vecindad del punto fijo, el m\'aximo de
$|g'|$ sobre todo el intervalo $[0.8,1.4]$ supera 1 en los extremos
(la ra\'iz de $g'$ en $x\to0.8$). Por eso se emplea la
\textbf{aceleraci\'on de Steffensen-Aitken}, que converge incluso cuando la
convergencia lineal ordinaria es lenta o marginalmente estable.

\subsection*{F\'ormula de Steffensen-Aitken ($\Delta^2$)}

Dado $x_n$, se calculan dos iteraciones adicionales:
\[
  x_n^{(1)} = g(x_n), \qquad x_n^{(2)} = g\!\left(x_n^{(1)}\right)
\]
y se extrapola:
\[
  \hat{x}_n = x_n - \frac{\bigl(x_n^{(1)}-x_n\bigr)^2}{x_n^{(2)}-2x_n^{(1)}+x_n}
\]

\subsection*{Tabla de iteraciones}

\begin{center}
\begin{tabular}{ccccc r}
\toprule
$n$ & $x_n$ & $g(x_n)$ & $g(g(x_n))$ & $\hat{x}_n$ & $|\hat{x}_n - x_n|$ \\
\midrule
1 & 1.000\,000\,00 & 0.960\,478\,66 & 0.950\,736\,42 & 0.947\,549\,25 & $5.25\times10^{-2}$ \\
2 & 0.948\,754\,93 & 0.948\,277\,00 & 0.948\,406\,07 & 0.948\,433\,89 & $8.85\times10^{-4}$ \\
3 & 0.948\,433\,89 & 0.948\,434\,04 & 0.948\,434\,06 & 0.948\,434\,07 & $1.77\times10^{-7}$ \\
\bottomrule
\end{tabular}
\end{center}

\[
  \boxed{x^* \approx 0.948\,434\,07, \qquad |f(x^*)| \approx 1.22\times10^{-14}}
\]

La convergencia se alcanz\'o en \textbf{3 iteraciones} con precisi\'on de
6 cifras significativas.

% ══════════════════════════════════════════════════════════════════════════════
\section*{Parte (b) --- Newton-Raphson}
% ══════════════════════════════════════════════════════════════════════════════

\subsection*{Derivada de $f$}

\[
  f'(x) = \frac{d}{dx}\bigl[2x\cos x - (x-2)^2\bigr]
        = 2\cos x - 2x\sin x - 2(x-2)
\]

\subsection*{Iteraci\'on}

\[
  x_{n+1} = x_n - \frac{f(x_n)}{f'(x_n)}
\]

\subsection*{Tabla de iteraciones ($x_0 = 1$, tol\,$= 10^{-8}$)}

\begin{center}
\renewcommand{\arraystretch}{1.25}
\begin{tabular}{c r r r r r}
\toprule
$n$ & $x_n$ & $f(x_n)$ & $f'(x_n)$ & $x_{n+1}$ & $|x_{n+1}-x_n|$ \\
\midrule
1 & 1.000\,000\,000\,0 & $+0.080\,605$ & $1.397\,663$ & 0.942\,328\,993\,2 & $5.77\times10^{-2}$ \\
2 & 0.942\,328\,993\,2 & $-0.010\,667$ & $1.766\,598$ & 0.948\,367\,107\,0 & $6.04\times10^{-3}$ \\
3 & 0.948\,367\,107\,0 & $-0.000\,116$ & $1.728\,257$ & 0.948\,434\,062\,0 & $6.70\times10^{-5}$ \\
4 & 0.948\,434\,062\,0 & $-1.2\times10^{-8}$ & $1.727\,832$ & 0.948\,434\,069\,9 & $8.25\times10^{-9}$ \\
\bottomrule
\end{tabular}
\end{center}

\[
  \boxed{x^* \approx 0.948\,434\,069\,9, \qquad |f(x^*)| = 0}
\]

Convergencia en \textbf{4 iteraciones} con tolerancia $<10^{-8}$.

% ══════════════════════════════════════════════════════════════════════════════
\section*{An\'alisis comparativo}
% ══════════════════════════════════════════════════════════════════════════════

\begin{center}
\renewcommand{\arraystretch}{1.3}
\begin{tabular}{lcc}
\toprule
\textbf{Criterio} & \textbf{Steffensen-Aitken} & \textbf{Newton-Raphson} \\
\midrule
Iteraciones hasta tol.\,$6\times10^{-7}$ & 3 & 3 \\
Iteraciones hasta tol.\,$10^{-8}$        & $\sim$4 & 4 \\
Evaluaciones de $f$ por iteraci\'on      & 2 ($g$ y $g\circ g$) & 2 ($f$ y $f'$) \\
Requiere derivada expl\'icita            & No & S\'i \\
Orden de convergencia                    & Superlineal ($\sim1.62$) & Cuadr\'atico ($\sim2$) \\
Riesgo de divergencia                    & Bajo (extra para aceleraci\'on) & Medio (si $f'\approx0$) \\
Dificultad de implementaci\'on           & Media & Baja \\
\bottomrule
\end{tabular}
\end{center}

\paragraph{Conclusi\'on.}
Ambos m\'etodos convergen a la misma ra\'iz
$x^* \approx 0.948\,434\,070$ con igual precisi\'on final.
\textbf{Newton-Raphson} es m\'as simple y tiene convergencia cuadr\'atica
te\'orica, pero requiere calcular $f'(x)$ anal\'iticamente.
\textbf{Steffensen-Aitken} evita la derivada expl\'icita y es robusto cuando
la contracci\'on de $g$ no es estricta en todo el intervalo, al precio de
requerir dos evaluaciones de $g$ por iteraci\'on en lugar de una.
Para esta funci\'on, la diferencia de eficiencia entre ambos es m\'inima.

% ══════════════════════════════════════════════════════════════════════════════
\section*{Ap\'endice: c\'odigo Python}
% ══════════════════════════════════════════════════════════════════════════════

\begin{lstlisting}[caption={Steffensen-Aitken y Newton-Raphson para f(x)}]
import math

def f(x):  return 2*x*math.cos(x) - (x - 2)**2
def df(x): return 2*math.cos(x) - 2*x*math.sin(x) - 2*(x - 2)
def g(x):  return 2 - math.sqrt(2*x*math.cos(x))

# Steffensen-Aitken
x = 1.0
for i in range(30):
    x0, x1, x2 = x, g(x), g(g(x))
    d = x2 - 2*x1 + x0
    if abs(d) < 1e-15: break
    xh = x0 - (x1 - x0)**2 / d
    if abs(xh - x) < 5e-7: break
    x = xh

# Newton-Raphson
x = 1.0
for _ in range(20):
    xn = x - f(x) / df(x)
    if abs(xn - x) < 1e-8: break
    x = xn
\end{lstlisting}

\end{document}
